\section{Introduction}
\label{sec:introduction}

eBPF programs now sit on hot paths for networking, observability, security, load balancing, and profiling. They run inside the Linux kernel, but they are not ordinary kernel code: each program is first checked by the kernel verifier, lowered by the in-kernel JIT, attached by an application-specific loader, and exercised by a workload whose behavior may be far from the microbenchmark that motivated the optimization. This stack creates a difficult optimization target. The same bytecode rewrite can be accepted or rejected by the verifier, neutral or profitable after JIT lowering, and useful or harmful depending on the workload's event mix.

Prior eBPF optimization systems show that this space has real headroom. K2 optimizes BPF bytecode before load~\cite{xu2021synthesizing}; Merlin optimizes across LLVM IR and BPF bytecode~\cite{mao2024merlin}; and safe-language or synthesis systems change how eBPF programs are produced~\cite{jia2025rex,zheng2024kgent}. These systems answer important questions about how to generate or optimize eBPF. They do not answer the question that motivates this paper: can an agent tune an already existing eBPF program when the only trustworthy feedback is the real verifier, the real JIT, and the real workload?

That question is attractive because eBPF optimization is full of local decisions. An optimizer must choose whether to run a pass at all, where to gate a rewrite, whether a kinsn-backed idiom is profitable on the target architecture, whether a verifier-state-dependent specialization is safe, and whether a measured speedup clears the noise floor. The search space is small enough to be actionable but irregular enough that a static all-apply policy is brittle. This is the shape of problem where agents look promising: they can inspect logs, form hypotheses, run constrained experiments, and update a policy from feedback.

However, evaluating such agents with ordinary code-generation benchmarks is the wrong abstraction. A generated eBPF program that compiles is not necessarily verifier-accepted. A verifier-accepted candidate can still break the application or the workload. A workload-correct candidate may only chase noise. Conversely, a candidate that looks useless by local rewrite count may improve a hot caller because tail-called descendants are accounted through the caller. Agent evaluation therefore needs a benchmark whose oracle composes verifier acceptance, application behavior, workload semantics, and performance measurement.

\sys is not a retitling of an existing benchmark runner. The existing runner is an execution substrate: it can build kernels and applications, run workloads, collect raw BPF counters, and preserve verifier/JIT/workload artifacts. An agent benchmark needs additional assets: frozen tasks, a base snapshot, instructions, an isolated evaluator the agent cannot edit, protected workloads and scoring scripts, hidden semantic checks, and traceable model/tool usage. \sys adds this missing layer for optimization tasks.

We propose \sys, a benchmark for agentic optimization of existing eBPF programs. \sys treats optimization as a closed-loop decision problem. The agent observes task context and prior feedback, proposes an allowed optimization action, the executor runs the real application and workload, and the oracle reports whether the candidate was accepted, correct, and faster under the benchmark's methodology. The first \sys track targets the layers already supported by our infrastructure: bytecode rewrites, live post-load ReJIT, and optional kinsn-backed JIT capabilities. The benchmark is not defined by those mechanisms. Source-level, LLVM-backend, and pre-load bytecode transforms are modeled as additional action adapters over the same oracle.

Our current data argues for this benchmark rather than for a premature optimizer claim. Across completed corpus runs, no standalone pass yet has a paper-ready measurable improvement over the measured noise floor. A noop ReJIT run and a skip-ReJIT floor differ by a suite-level amplitude of 0.0431 in the post/baseline ratio, while some app-level cells vary much more. Applied-site counts also fail as a profitability proxy: a full-corpus run was near flat overall, yet changed-code rows regressed while unchanged rows improved, revealing residual noise and phase effects. Policy tuning helps on selected dense benchmarks, but the corpus remains mixed. These results say that eBPF optimization is exactly the kind of constrained, feedback-driven problem a benchmark should expose.

\noindent\textbf{Contributions.} This paper makes four contributions:
\begin{itemize}
  \item We define eBPF optimization for agents as a layered action space spanning source, LLVM backend, pre-load bytecode, live bytecode/ReJIT, and kernel/JIT/kinsn transformations.
  \item We present \sys, a benchmark interface with task manifests, allowed actions, hidden evaluator checks, feedback traces, and a composed oracle over verifier, application, workload, and performance outcomes.
  \item We use existing corpus results to motivate the benchmark, showing that static pass policies, rewrite counts, and verifier acceptance alone are insufficient evidence of optimization success.
  \item We outline an evaluation plan and scoreboard for comparing no-op, static, random, human-tuned, and LLM-agent policies under equal run budgets.
\end{itemize}
